\answer{ex:13:1}{$q=100^{1/99}\approx1.048$, шаг $\approx4.8\%$; диапазон $\approx2.2\cdot10^3$ ($67$~дБ); $\approx15$ включений на удвоение; у линейного ЦАП шаг $22f_1$, т.\,е. $220\%$ при $10f_1$; точность $10\%$ достигается при $n\ge53$.}
\answer{ex:13:2}{$P_{\text{сбоя}}\approx0.47$, $1.8\cdot10^{-4}$, $2.8\cdot10^{-16}$; при $S=2$ около $300$ сбоев в сутки, при $S=4$ около $5\cdot10^{-10}$ в сутки; оценка Чернова даёт $\le10^{-14}$.}
\answer{ex:13:3}{$H(s)=eF_1T_c/(1+sT_c)^2$, частота среза $1/(2\pi T_c)\approx3.2$~Гц; $\bar F=e\,rT_cF_1$ ($2.0F_1$ и $5.4F_1$); при $5$~имп/с $F$ от $0.21F_1$ до $1.10F_1$; размах $10\%$ при $r\approx22$~имп/с (по первой гармонике $\approx20$).}
\answer{ex:13:4}{$Gd<\pi/2$, $f=1/(4d)$; быстрее всего без колебаний при $Gd=1/e$, скорость затухания $1/d$; $d=30\ms$: $G\approx12$~с$^{-1}$, $\tau=30\ms$; $d=150\ms$: $G\approx2.5$~с$^{-1}$, $\tau=150\ms$.}
\answer{ex:13:5}{$E\approx0.427$, $T\approx1.70\,\tau_a=0.68$~с (модель с $\tau=20\ms$ даёт $0.84$~с); ритм существует при $b>\beta-1$; $T$ не зависит от $I$.}
\answer{ex:13:6}{$T\approx0.48$, $0.84$, $1.55$~с при $\tau_a=0.2$, $0.4$, $0.8$~с ($0.39$, $0.73$, $1.42$~с при $\tau=5\ms$); при любом $I$ $T=0.840$~с; при $b=0.8$ ритма нет, при $b=1.05$, $1.2$, $1.5$, $2$, $4$ $T\approx2.73$, $1.74$, $1.19$, $0.84$, $0.45$~с.}
\answer{ex:13:7}{$a\ge33.3$~с$^{-1}$, $F_1\approx58$~Н, $F_2\approx8.3$~Н, $G=e^{-2}/d\approx4.5$~с$^{-1}$; без мышцы-противника $\tau\approx17\ms$, требование не выполнено; предел $\tau\approx9.2\ms$; чем сильнее совместное напряжение, тем меньше надо брать $G$.}
\answer{ex:13:8}{Открытая задача. Например, порог в накачке усталости $b[r_i-r_0]_+$ даёт $I_{\min}=\beta b r_0/(b-\beta+1)\approx0.53$ при $b=4$, $r_0=0.2$ и период от $1.8$~с при $I=0.6$ до $0.52$~с при $I=3$.}
